\section{What the Approval Obliges---and What It Does Not}

\subsection{Public Revelation is complete}

The doctrinal frame is prior to every fact in this study and settles
more of the La Salette question than any archival discovery could.
Public Revelation is complete in Christ: ``the Christian economy,
since it is the new and definitive covenant, will never pass away;
and no new public revelation is to be expected before the glorious
manifestation of our Lord, Jesus Christ.'' Private revelations, even
recognized ones, do not belong to the deposit of faith; their purpose
is not ``to improve or complete Christ's definitive Revelation, but
to help live more fully by it in a certain period of
history.''\endnote{\work{Dei Verbum} 4, quoted at n.~1 of the 2024
norms, which cite \work{AAS} 58 (1966), p.~819; \work{Catechism of
the Catholic Church} n.~67, quoted at n.~4 of the same norms. Both
are read here in the Holy See's English text of the norms
(References), read 2026-07-25.}

Everything at La Salette---the 1846 discourse, the 1851 manuscripts,
the 1879 booklet, every reported locution of either seer for the
remainder of their lives---falls on the private-revelation side of
that line. Nothing there can add an article to the faith, correct
the Gospel, or bind a conscience as Revelation binds it.

\subsection{\term{Fides humana}: the Holy See's own word for La
Salette}

Most apparition studies must reason their way to the assent
question. Here the Holy See supplied the answer in 1877 in a text
that names La Salette explicitly (\S8.3): such apparitions or
revelations have been neither approved nor reprobated by the
Apostolic See, ``but only permitted as to be piously believed
\term{fide solum humana}''---by merely human faith. That is the
technical term, in a Roman act, about this event.

The modern formulation is the same doctrine in modern words. The
2024 norms state as a rule that ecclesiastical authority does not
give ``a positive recognition of the divine origin'' of such
phenomena (n.~11); that even where a \term{Nihil obstat} is granted
``such phenomena do not become objects of faith, which means the
faithful are not obliged to give an assent of faith to them''
(n.~12); and, quoting Benedict~XVI, that approval ``essentially
means that its message contains nothing contrary to faith and
morals; it is licit to make it public and the faithful are
authorized to give to it their adherence in a prudent manner\ldots\
It is a help which is proffered, but its use is not
obligatory.''\endnote{2024 norms, nn.~11--12 and note~24, quoting
Benedict~XVI, \work{Verbum Domini} (30 September 2010), n.~14:
\work{AAS} 102 (2010), p.~696; the \work{Verbum Domini} paragraph is
quoted in full in the norms' own footnote and is taken from there,
not from an independent reading of the exhortation (Appendix~A).}

Against that background the 1851 formula reads exactly as the
Dicastery says it reads. \term{Les fidèles sont fondés à la croire
indubitable et certaine} states that the faithful have \emph{grounds}
for a belief; the grammar of the sentence is a warrant, not a
precept. The Dicastery's 2024 presentation nevertheless names this
very sentence as an example of episcopal language that ``conflicted
with the Church's own conviction that the faithful did not have to
accept the authenticity of these events''---which is a candid
admission that formulas of that period could be, and were, heard as
obligations. Both things are true, and this study says both: the act
warrants belief, and the Church now judges that its wording
overreached what it could ask.\endnote{2024 norms, presentation,
``Reasons for the New Norms.'' It is worth noting precisely what the
Dicastery does and does not do with the 1851 sentence: it cites it
as an example of a style of statement, in a historical explanation of
why new norms were needed. It does not annul the 1851 act, reopen the
case, or assign the event a new determination, and this study asserts
no such effect.}

Concretely: a Catholic is free---reasonably, prudently, without
disloyalty---to believe that the Blessed Virgin appeared above La
Salette on 19 September 1846, and free not to. What no Catholic is
free to do is invert the order of assent, treating any La Salette
text as a second deposit of faith, an oracle settling doctrine, a
prediction binding the future, or a warrant for judging the Church's
pastors.

\subsection{The exact object of each standing act}

The inventory below is the whole dossier read in one column: every
act this study located, with the object that act itself establishes
and the witness in which this study read it. Nothing is inferred
from proximity, and no act in it carries any object but its own.

\begin{actstable}
1851, 19 Sept. & Bishop of Grenoble, doctrinal mandement & That the
apparition of 19 September 1846 to two shepherd children bears all
the marks of truth and that the faithful have grounds for believing
it indubitable and certain; cult authorized; devotional publications
reserved; public contest forbidden in the diocese & Sanctuary
transcription; DDF 2024 citation; Ullathorne 1854; Rousselot 1854
(It.) \\
1852, 1 May & Bishop of Grenoble, mandement & The Missionaries of
Our Lady of La Salette instituted; the sanctuary and hospice
undertaken & Institute's published text \\
1852, Aug.--Dec. & Pius~IX, rescripts, briefs, indult & Privileged
altar, votive Mass, indulgences, faculties, archconfraternity, and
the 19 September diocesan solemnization: acts of cult, of their own
period's discipline & \work{Vie de M.~Rousselot} (1866) \\
1877, 12 May & Sacred Congregation of Rites & That the Apostolic See
had neither approved nor condemned the apparitions, but permitted
them as piously believable by merely human faith; images and
sodalities permitted; a first-class feast with octave and special
litanies of the apparition refused & \work{ASS} 11 (1878) 512--514
\\
1915, 21 Dec. & Holy Office, decree & A universal precept forbidding
the faithful to discuss or treat of the \term{Secret de la Salette}
in publication, with penalties; devotion under the title
\term{Reconciliatrix} expressly preserved & \work{AAS} 7 (1915) 594
\\
1916, 12 Apr.\slash 5 June & Holy Office; Congregation of the Index
& A named two-volume exegesis of the secret condemned and placed on
the Index & \work{AAS} 8 (1916) 175, 178--179 \\
1923, 9 May & Holy Office, decree & A named 1922 reprint of
Mélanie's text proscribed and condemned, with an order to withdraw
copies from the faithful & \work{AAS} 15 (1923) 287--288 \\
1945, 8 Oct. & Pius~XII, letter & Papal association with the
centenary; the 1851 act described as a diocesan canonical process
that proved favourable; the event kept in the reported register &
\work{AAS} 38 (1946) 155--156 \\
1961, 16 June & John~XXIII, apostolic letter & Our Lady of La
Salette, invoked as Reconciler of sinners, constituted principal
patroness of one religious congregation & \work{AAS} 53 (1961)
606--607 \\
1996, 6 May and 9 July & John Paul~II, letters & Papal association
with the 150th anniversary; a legate appointed; the event kept in
the reported register (\term{ab Ecclesia Matre perhibetur}) & Holy
See web texts \\
2016, 18 Mar. & Congregation for Divine Worship & The celebration
inscribed in the Proper of France on 19 September as an optional
memorial & French episcopal conference's site \\
\end{actstable}

One item of \S5 is deliberately absent from that table: the letter
of 14 August 1880 (\S5.6). It is not a standing act---no text of it
was reached, it was never published in any gazette, and it is known
only as a partisan witness of 1882 reports it---and a table of
located acts is not the place for it.

Read down that column of objects and the shape of the dossier is
plain. Not one act in it declares the apparition supernatural in the
sense the 2024 norms describe; not one approves a secret; not one
obliges belief; and the only universal acts the Holy See ever issued
on the subject were restrictive of the secret literature and
protective of the devotion.

\subsection{Devotional practice within the boundary}

What the approved devotion asks can be stated entirely inside
ordinary Catholic doctrine, and the custodians state it that way
themselves. The message's own asks are Sunday, reverence for the
name of Christ, daily prayer, penance, and Mass---the substance of
the third and second commandments and of the Church's ordinary
precepts. The shrine's own charter defines its priests by the
ministry of reconciliation and the Eucharist (\S8.1). The title under
which the Church honours the event is \term{Réconciliatrice}, read
within Christ's unique mediation (\S3.6). John Paul~II's
sesquicentenary letter reads the whole event through the crucifix on
the figure's breast (\S9.5). There is no devotional practice proper
to La Salette that stands apart from Sunday, confession, the
Eucharist, and prayer.

Two boundaries deserve their own sentences. First, weeping is not a
doctrine. The tears reported at La Salette are, in John Paul~II's
reading, a pedagogy of the gravity of sin and of the Son's wounded
fidelity; they are not evidence of a divided will between the Father,
the Son, and the Mother, and any presentation of La Salette in which
Mary restrains an angry Christ from a reluctant Father has left
Catholic doctrine, whatever imagery it borrows from the discourse.
Second, threatened chastisement is conditional, and it is
conditional in the discourse's own grammar (``if my people will not
submit''; ``if they are converted''). Using La Salette to date, to
predict, to alarm, or to identify culprits is using it against its
own text.\endnote{Editorial synthesis, labelled as such, over the
elements cited at \S3.5, \S3.6, \S8.1, and \S9.5.
\work{Lumen gentium} 60--62 and CCC 970 supply the doctrinal control
on the Marian imagery; nothing in the \emph{public} La Salette
message strains them, which is a statement about the public message
only and is not extended to any other text.}

\subsection{Disagreement preserved}

A judgment-first study must say where the record leaves room for
honest disagreement, and here it leaves a good deal.

The 1851 act argued from evidence---the children's constancy, the
concourse, the reported prodigies---and evidence-based judgments can
be weighed differently by serious minds; the Church has never
required anyone to find those arguments conclusive, and in 1877 the
Holy See declined to call the event approved at all. The seers'
later lives are genuinely contested and the sources are polarized
(\S7). The relation between the 1851 manuscripts and the 1879 text
is contested, and the promoter's own admissions make it so (\S5.4).
The present force of the 1915 decree is unsettled (\S5.9). And this
study's own reach is limited in ways Appendix~A sets out---most
importantly, no French printing of the integral mandement and no
access to the primary documentary edition of the early records.

What the Church's acts secure is narrower and firmer than victory in
any of those disputes: that a Catholic may, with the Church,
prudently believe that the Blessed Virgin appeared at La Salette;
that the devotion under the title of Reconciler is safe and
protected; that the shrine and its ministry of reconciliation are
good; and that nobody is obliged to believe anything about a secret.
Nothing in this study should be read as claiming more.
